\documentclass[11pt]{article}
\usepackage[a4paper,margin=26mm]{geometry}
\usepackage[T1]{fontenc}
\usepackage[utf8]{inputenc}
\usepackage{lmodern,amsmath,amssymb,amsthm,mathtools,microtype,booktabs,array}
\usepackage{hyperref}
\hypersetup{colorlinks=true,linkcolor=blue,citecolor=blue,urlcolor=blue}
\newtheorem{theorem}{Theorem}
\newtheorem{lemma}[theorem]{Lemma}
\newtheorem{proposition}[theorem]{Proposition}
\newtheorem{corollary}[theorem]{Corollary}
\theoremstyle{remark}\newtheorem{remark}[theorem]{Remark}
\DeclareMathOperator{\diag}{diag}
\DeclareMathOperator{\Hess}{Hess}
\DeclareMathOperator{\rank}{rank}
\newcommand{\R}{\mathbb R}
\newcommand{\Q}{\mathbb Q}
\newcommand{\one}{\mathbf 1}
\title{Uniform ten-spin Lorentzian obstructions\\for cycle clique blowups}
\author{Lluis Eriksson\\\small Independent researcher, Sweden}
\date{17 October 2026\\\small Working-paper manuscript}
\begin{document}
\maketitle
\begin{abstract}
For every $m\ge4$ and $s\ge2$, we construct a strictly positive
rational antiferromagnetic target of rank two on at most ten states
whose partition polynomial for $C_m[K_s]$ is not Lorentzian.
The state bound is independent of both parameters; seven states
suffice when $s\ge9$, and six suffice for even $m$.
The new case is an odd cycle, whose edges are critical for its
independence number and hence fall outside the earlier
independence-stable blowup theorem.
We classify the minimum-induced-edge occupations and compress
their leading Hessian kernel to an initial block of at most eight
profiles. For $s\ge20$, a five-profile block has inertia $(2,3,0)$,
proved by six exact coefficient certificates in the cycle length,
clique size and central binomial coefficient. Eighteen bounded
clique sizes are covered by polynomial certificates for infinite
cycle-length tails and 42 explicit integer exceptions.
We retain failures of smaller blocks and give a positive seven-state
witness for $C_5[K_9]$. The target may depend on $m,s$; neither
optimality nor the general connected-source characterization is claimed.
\end{abstract}

\section{The result and its scope}
For a finite simple graph $G=(V,E)$ and a symmetric nonnegative
$q\times q$ matrix $K$, write
\begin{equation}
 Z_G(\mathbf x;K)=
 \sum_{\sigma:V\to\{0,\ldots,q-1\}}
 \prod_{uv\in E}K_{\sigma(u),\sigma(v)}
 \prod_{v\in V}x_{\sigma(v)}.                         \label{eq:Z}
\end{equation}
We call $K$ antiferromagnetic when it has at most one positive
eigenvalue. Inertia is written in the order positive, negative, zero.
The graph $C_m[K_s]$ replaces each vertex of a simple cycle by a
clique of order $s$, joining all pairs across adjacent classes.

For a homogeneous polynomial with strictly positive coefficients
on the full degree simplex, Lorentzianity is equivalent to every
quadratic partial derivative having a Hessian with at most one
positive eigenvalue \cite[Theorem 2.25]{BH}. Positive targets in
\eqref{eq:Z} give this full support. Lee--Oh--Seo proved that clique
sources are Lorentzian for every antiferromagnetic target
\cite[Theorem 3.1]{LOS}. The connected-source converse discussed
in \cite[Section 6]{LS} asks whether only cliques have this property.
We address one two-parameter family.

\begin{theorem}\label{thm:main}
Let $m\ge4$ and $s\ge2$ be integers. There are $q\le10$ distinct
rational nodes $a_0,\ldots,a_{q-1}\in(-1,1)$, one of them zero,
such that
\[
 K=\one\one^\top-aa^\top
\]
has inertia $(1,1,q-2)$ and $Z_{C_m[K_s]}(\mathbf x;K)$ is not
Lorentzian. One may take $q\le7$ if $s\ge9$ and $q=6$ if $m$ is even.
The matrix is strictly positive entrywise.
\end{theorem}

The parameters of $K$ may depend on both $m$ and $s$.
Adding distinct nodes in $(-1,1)$ allows exactly ten states:
the restriction to the old variables preserves the old Hessian
obstruction. Combining the theorem with the clique result gives
a characterization within this family.

\begin{corollary}
For $m\ge3$ and $s\ge2$, the polynomial
$Z_{C_m[K_s]}(\mathbf x;K)$ is Lorentzian for every antiferromagnetic
target $K$ if and only if $m=3$.
\end{corollary}
\begin{proof}
For $m=3$ the source is $K_{3s}$, so \cite[Theorem 3.1]{LOS}
applies. Theorem~\ref{thm:main} supplies the converse.
\end{proof}

The earlier conditioned-profile method \cite{RankTwo} and its
six-spin refinement \cite{SixSpin} are inputs, not new claims here.
The latter covers bases $F$ for which
$\alpha(F-e)=\alpha(F)$ for every edge.
Even cycles satisfy this hypothesis. An odd cycle $C_{2k+1}$ does
not: deleting any edge gives a path with independence number
$k+1$, whereas $\alpha(C_{2k+1})=k$.
Thus the odd case requires a different ground-layer analysis.
Each nontrivial clique class also contains maximum-degree true
twins, so the earlier no-maximum-degree-twins criterion does not
settle this family.

The contribution is the uniform odd-cycle construction and the
bound on its number of states. Turan equality, generalized
Vandermonde invertibility and continuity of inertia are classical.
The homomorphism inequalities for clique blowups in \cite{LOS,LS}
concern a different assertion and are not contradicted by a failure
of Lorentzianity of the full source partition polynomial.
Our checks of these sources and the preceding manuscripts support
this comparison; they do not establish exhaustive worldwide priority.

\section{From an initial profile block to a small target}
We give the transfer argument for an arbitrary number of active
profiles, including the order of its two limiting operations.

Let $n=|V(G)|$, $0\le h\le n-2$, $a_0=\ell\ne0$ and $a_1=0$.
Put $\theta=\ell^2$. For $S\subseteq V(G)$ of size $h$, let
$r_S(v)=|N(v)\cap S|$ for $v\notin S$.
Let $\nu^S_r$ count such vertices of profile $r$, and let $C^S_{rt}$
count ordered adjacent pairs of distinct outside vertices with
profiles $r,t$. Pad $\nu^S$ by a zero in coordinate $h+1$.
The matrix $D$ of size $(h+2)\times(h+1)$ is given by
$De_r=e_r-e_{r+1}$. Define
\begin{equation}
 Q_h(\theta)=\sum_{|S|=h}(1-\theta)^{e(G[S])}
 \left[\theta\{\nu^S(\nu^S)^\top-\diag\nu^S\}
             -DC^SD^\top\right].                       \label{eq:Q}
\end{equation}
With $U_{ir}=(1-\ell a_i)^r$, direct differentiation gives
\begin{equation}
 \Hess\!\left(\partial_0^h\partial_1^{n-h-2}Z_G\right)
 =\frac{h!(n-h-2)!}{\theta}\,UQ_h(\theta)U^\top.          \label{eq:H}
\end{equation}
Indeed, fixing the $h$ vertices assigned to the pin spin contributes
$(1-\theta)^{e(G[S])}$. The two undifferentiated vertices give
the profile products, and $\nu\nu^\top-\diag\nu$ counts all ordered
distinct pairs. The edge correction is $a_i a_j$ times the ordered
edge count; the identity $(UD)_{ir}=\ell a_i(1-\ell a_i)^r$
produces the second term. Edges involving vertices differentiated
in the zero-node spin have weight one. This also accounts for
the factorials, including coincident residual spins.

Set $e_0=\min_{|S|=h}e(G[S])$. The leading kernel is
\begin{equation}
 L=\lim_{\theta\uparrow1}\frac{Q_h(\theta)}{(1-\theta)^{e_0}}
 =\sum_{\substack{|S|=h\\e(G[S])=e_0}}
   \{\nu^S(\nu^S)^\top-\diag\nu^S-DC^SD^\top\}.           \label{eq:L}
\end{equation}
An index is \emph{active} if its row in $L$ is nonzero.
In particular, activity is not defined just by $\nu_r>0$.

\begin{lemma}[Initial-profile compression]\label{lem:compression}
Suppose the principal submatrix of $L$ on its first $r$ active
indices is nonsingular and has at least two positive eigenvalues.
Then $G$ has a strictly positive rational rank-two
antiferromagnetic obstruction on $r+2$ states.
\end{lemma}
\begin{proof}
Let the first active indices be $b_1<\cdots<b_r$. For rational
$\delta>0$, use the $r$ auxiliary nodes $a_{i+1}=-1+\delta i$,
$i=1,\ldots,r$. At $\ell=-1$ their profile evaluation matrix
has entries $(\delta i)^b$. Its columns on the selected indices
form
\[
 V_\delta=W\diag(\delta^{b_1},\ldots,\delta^{b_r}),
 \qquad W_{ij}=i^{b_j}.
\]
The matrix $W$ is invertible: a nonzero polynomial with at most
$r$ monomials has at most $r-1$ distinct positive roots, by
Descartes' rule of signs. For every omitted active index $b>b_r$,
the $j$th entry of $V_\delta^{-1}((\delta i)^b)_i$ is a constant
times $\delta^{b-b_j}$, and therefore tends to zero.
Inactive rows contribute nothing. Congruencing the auxiliary
restriction of $ULU^\top$ by $V_\delta^{-1}$ thus gives the selected
block in the limit $\delta\downarrow0$.

Choose rational $\delta$, with $0<\delta<1/r$, sufficiently small
to retain its inertia. Now fix $\delta$ and put
$\ell=-1+\varepsilon$ with $\varepsilon\downarrow0$.
The positive normalization in \eqref{eq:H} and \eqref{eq:L}
and continuity retain at least two positive eigenvalues for
rational $\ell$ sufficiently close to $-1$.
We also impose $0<\ell+1<\delta$, so the pin, zero node and
auxiliary nodes are distinct. All nodes lie in $(-1,1)$, hence
$K_{ij}=1-a_i a_j>0$.
The two vectors $\one,a$ are linearly independent; Sylvester's
law of inertia gives $(1,1,r)$ for $K$.
The offending principal Hessian restriction proves non-Lorentzianity.
\end{proof}
No evaluation at a zero-entry boundary target is asserted.
The target construction is terminating in principle by rational
search and exact inertia, but no useful denominator bound is proved.

\section{Minimum-edge occupations on an odd cycle}
For the rest of the proof put $m=2k+1$, $k\ge2$, and
\begin{equation}
 n=(2k+1)s,\qquad h=ks-1,\qquad
 e_0=\frac{(s-1)(ks-2)}2.                              \label{eq:parameters}
\end{equation}
Both derivative orders $h$ and $n-h-2=(k+1)s-1$ are nonnegative.

An occupation vector $c=(c_0,\ldots,c_{m-1})$ records how many
vertices of each clique class belong to $S$. Its weight, the number
of labelled subsets it represents, is $\prod_i\binom{s}{c_i}$.

\begin{lemma}\label{lem:ground}
The minimum number of edges in an $h$-vertex induced subgraph of
$C_m[K_s]$ is \eqref{eq:parameters}. Every minimizing subgraph is
a union of $k$ clique components, one of size $s-1$ and the others
of size $s$. Its occupation is either
\begin{enumerate}
\item one component occupying two consecutive cycle classes, with
all other components occupying one class and all intervening gaps
having length one; or
\item every component occupying one class, with one gap of length
two and all other gaps of length one.
\end{enumerate}
\end{lemma}
\begin{proof}
The independence number of the source is $k$.
Turan's theorem applied to the complement of an induced
$h$-vertex graph gives the minimum in \eqref{eq:parameters};
equality requires $k$ balanced clique components \cite{Turan}.
Selecting full classes of a maximum independent set except for
one omitted vertex attains equality.
Since $m\ge5$, a clique component can use at most two consecutive
cycle classes. Different components require at least one empty
class between them. There are $2k+1$ classes and $k$ components,
so at most one component can occupy two classes. If there is
one, all $k$ gaps have length one. If there is none, precisely
one gap has length two. These conditions are also sufficient.
\end{proof}

Here is an exact summation with convenient endpoint corrections.
Anchor a possible split at classes $0,1$ and place the other
components at $3,5,\ldots,2k-1$.
\begin{itemize}
\item Type A has occupations $(t,s-1-t)$ at $0,1$ and all other
components full, with $0\le t\le s-1$. Its weight is
$\binom{s}{t}\binom{s}{t+1}$.
\item Type B has occupations $(t,s-t)$ at $0,1$, with
$0\le t\le s$, and one of the $k-1$ other components partial
of size $s-1$. Its weight is $s\binom{s}{t}^2$.
\end{itemize}
Take all $m$ rotations. Strict splits occur exactly once.
Each unsplit configuration occurs twice among the endpoints:
one can regard either class next to its unique double gap as
the degenerate split. Subtract one copy of all unsplit configurations.
They are the $m$ rotations of the $k$ choices of a partial class
in $\{0,2,\ldots,2k-2\}$, each with weight $s$.
Thus the number of minimizing labelled subsets is
\begin{equation}
 N=m\left[\binom{2s}{s-1}+(k-1)s\binom{2s}{s}-ks\right]. \label{eq:N}
\end{equation}
The signed representation counts a positive combinatorial set
exactly once; it does not assign negative physical weights.

For an occupation define the remaining class size and its profile by
\[
 b_i=s-c_i,\qquad \rho_i=c_{i-1}+c_i+c_{i+1}.
\]
Indices are cyclic. With zero padding as needed, put
\[
 \nu=\sum_i b_i e_{\rho_i},\qquad
 C=\sum_i b_i(b_i-1)e_{\rho_i}e_{\rho_i}^{\top}
 +\sum_i b_i b_{i+1}
 (e_{\rho_i}e_{\rho_{i+1}}^\top+e_{\rho_{i+1}}e_{\rho_i}^\top).
\]
The occupation contributes its weight times
\begin{equation}
 M(c)=\nu\nu^\top-\diag\nu- DCD^\top.                  \label{eq:occupation}
\end{equation}
The internal edges of an empty clique class contribute
$s(s-1)$ ordered pairs and must be retained.
Summing \eqref{eq:occupation} over the above types and corrections
is an explicit formula for $L$.
All active indices lie between $s-1$ and $2s+1$:
the remaining part of a selected component has profile $s-1$ or
$s$, an empty class has profile between $s-1$ and $2s$, and
$D$ shifts upwards by at most one.

\section{Five initial profiles for all large clique sizes}
Let $B=\binom{2s}{s}$ and
\[
 A=(s-1,s,s+1,s+2,s+3),\qquad
 T=L[A,A]/m.
\]
We describe a finite symbolic construction of $T$, rather than
fitting its entries from numerical examples.
For a pattern $c$, let $P(c)$ be the restriction of
\eqref{eq:occupation} to $A$, with
\begin{equation}
\begin{split}
 P(c)=\nu_A\nu_A^\top-\diag\nu_A
 -\sum_i b_i(b_i-1)d_i d_i^\top
 \\{}-\sum_i b_i b_{i+1}(d_i d_{i+1}^\top+d_{i+1}d_i^\top),
 \qquad d_i=(e_{\rho_i}-e_{\rho_i+1})|_A.
\end{split}                                           \label{eq:P}
\end{equation}

For a type A pattern, let $P_A^\circ$ be its generic projected
matrix when the two variable boundary profiles lie outside $A$.
The sum over all split parameters is
\begin{align}
 \mathcal A={sB\over s+1}P_A^\circ+
 \sum_{j=0}^{4}\binom{s}{j}\binom{s}{j+1}
 \bigl(P_A(j)+P_A(s-1-j)-2P_A^\circ\bigr).             \label{eq:Asum}
\end{align}
For a type B partial-position category $p$, the analogous sum is
\begin{equation}
 \mathcal B_p=sB P_{B,p}^\circ+
 \sum_{j=0}^{4}s\binom{s}{j}^2
 \bigl(P_{B,p}(j)+P_{B,p}(s-j)-2P_{B,p}^\circ\bigr).    \label{eq:Bsum}
\end{equation}
These formulas hold for $s\ge11$.
For example, in type A the variable profiles are $s+t$ and
$2s-1-t$. In type B they are $s+t$ and $2s-1-t$ for a nearby
partial class, and $s+t$ and $2s-t$ for a distant partial class.
For $k=2$ the unique partial class touches both boundaries,
giving $s+t-1$ and $2s-1-t$ instead.
Solving for when either $\rho$ or $\rho+1$ enters $A$ puts all
exceptions in the two endpoint lists in \eqref{eq:Asum}--\eqref{eq:Bsum}.
For $s\ge11$ these lists are disjoint. The weights outside
them sum by Vandermonde's identity; the correction sums add
the exact weighted difference from the generic matrix.

For $k\ge3$ there are two near partial positions and $k-3$ far
positions. The two near positions have equal \emph{summed} matrices
by reflection and $t\leftrightarrow s-t$.
For unsplit endpoint patterns there are two near partial positions
and $k-2$ far positions. Therefore
\begin{equation}
 T=\mathcal A+2\mathcal B_{\rm near}
 +(k-3)\mathcal B_{\rm far}
 -2sP_{E,\rm near}-(k-2)sP_{E,\rm far}.                \label{eq:Tlarge}
\end{equation}
Representatives may be drawn on $C_9$. Extra full components
contribute only profile $2s$, outside the projection for $s\ge11$.
The local near pattern is valid already at $k=3$; the absent
far B category has multiplier zero.
For $k=2$, compute the A pattern and the single B position on
$C_5$ and subtract its two distinct unsplit endpoint patterns.
It is not obtained by substituting $k=2$ into \eqref{eq:Tlarge}.

This defines every entry of the two rational-polynomial matrices
in $s,k,B$. Their expanded entries and determinant factorizations
are in \texttt{odd-initial-five-symbolic.json}.
The first row in both branches, with the appropriate value of $k$, is
\begin{equation}
 \left(0,\ Bks^2,\ 2ks^4,\
 {ks^4(s-1)^2\over2},\
 {ks^4(s-1)^2(s-2)^2\over18}\right).                  \label{eq:firstrow}
\end{equation}
All four displayed couplings are positive in the domain used below.
There are no lower profiles, so $A$ consists of the first five
active indices.

\begin{proposition}\label{prop:large}
For every integer $s\ge20$ and $k\ge2$, $T$ has inertia $(2,3,0)$.
\end{proposition}
\begin{proof}
Let $\Delta_j$ be its leading $j\times j$ determinant.
Exact expansion of \eqref{eq:P}--\eqref{eq:Tlarge} gives
\[
 \Delta_2=-B^2k^2s^4<0
\]
in both branches, including $k=2$.
For $\Delta_3,\Delta_4,\Delta_5$ the required signs are $+,-,-$.
Here is the exact finite certificate for those signs.
Take the numerator of the signed determinant, retaining its
positive constant denominator separately. Factor out only powers
of $s,s-1,s-2,k,B$, which are positive in the domain. Denote the
remaining signed polynomial by $F$. Substitute
\begin{equation}
 B=s^8+w,\qquad k=3+v,\qquad s=20+u,                 \label{eq:shift}
\end{equation}
omitting $v$ in the $k=2$ branch.
Every nonzero coefficient of the resulting six polynomials is
positive, and every constant term is positive. The complete
integer coefficient lists are provided, with the following counts.
\begin{center}
\begin{tabular}{@{}lrrr@{}}
\toprule
Branch & $\Delta_3$ & $-\Delta_4$ & $-\Delta_5$\\
\midrule
$k=2$ & 56 & 64 & 80\\
$k\ge3$ & 112 & 196 & 328\\
\bottomrule
\end{tabular}
\end{center}
These are 836 retained coefficients, not sampled values or rounded
eigenvalues. The supplement reconstructs the matrices,
determinants, stripped factors and all coefficients using exact
integer and rational arithmetic.

It remains to justify $w>0$ in \eqref{eq:shift}.
At $s=20$,
\[
 \binom{40}{20}-20^8=112246528820>0.
\]
The polynomial $(4s+2)s^8-(s+1)^9$ has strictly positive
coefficients in $s-20$, as the exact expansion in the supplement
shows. The recurrence $B(s+1)/B(s)=(4s+2)/(s+1)$ then proves
$B(s)>s^8$ by induction. Thus all six signed determinants are
strictly positive for the claimed integer parameters.

The initial $2\times2$ block is nonsingular with inertia $(1,1,0)$.
The subsequent Schur pivots have signs
$\Delta_3/\Delta_2<0$, $\Delta_4/\Delta_3<0$ and
$\Delta_5/\Delta_4>0$, proving $(2,3,0)$.
We do not use the zero first diagonal entry as a scalar pivot.
\end{proof}

\section{Bounded clique sizes and infinite cycle tails}
For $2\le s\le19$, a separate exact polynomial certificate
avoids the overlapping endpoints of the large-$s$ formula.
For fixed $s$, the full matrix $L/m$ is polynomial in $k$ for
$k\ge3$. To construct it, use the same A, B and endpoint categories
as before, now retaining \emph{all} profile coordinates.
A representative on $C_9$ is extended symbolically by
\begin{equation}
 \nu_{2s}\longmapsto\nu_{2s}+(k-4)s,\qquad
 C_{2s,2s}\longmapsto C_{2s,2s}+(k-4)s(s-1).           \label{eq:extra}
\end{equation}
An additional full component leaves no outside vertices in its
class; its following empty class contributes precisely the
vertices and internal ordered edges in \eqref{eq:extra}.
No other remaining adjacency or local profile changes.
The multipliers of the near and far B categories are $2$ and
$k-3$; those of the endpoint corrections are $2$ and $k-2$.
Reflection identifies the near categories only after summing
over $t$. At $k=3$ the far B term is absent; the other local
patterns have the same formula with one fewer ordinary empty class.
This is a combinatorial derivation of a polynomial matrix,
not interpolation. The case $k=2$ is evaluated directly.

For each $s$, take its first active $r$ coordinates.
The following table gives a certified block for the infinite
tail $k\ge N$, the largest block needed in the finitely many
integers $3\le k<N$, and the size for $k=2$.
\begin{center}
\begin{tabular}{@{}rrrrr@{}}
\toprule
$s$ & Tail size $r$ & $N$ & Largest finite size & $k=2$ size\\
\midrule
2 &4&3&--&4\\
3 &5&3&--&4\\
4 &5&8&6&5\\
5 &6&3&--&5\\
6 &7&3&--&6\\
7 &6&40&8&6\\
8 &6&3&--&6\\
9--19 &5&3&--&5\\
\bottomrule
\end{tabular}                                                     \label{tab:small}
\end{center}

For each row, the determinant of the indicated initial block
is a nonzero polynomial $d_s(k)$.
Its leading sign $\eta_s\in\{-1,1\}$ is chosen so that every
coefficient of $\eta_s d_s(N+u)$ is nonnegative and its constant
coefficient is positive. These complete coefficient lists are in
\texttt{odd-initial-small-certificates.json}.
Hence the block is nonsingular throughout the real interval
$[N,\infty)$. Exact symmetric elimination at $N$ gives at least
two positive eigenvalues. Since inertia cannot change along a
continuous nonsingular real symmetric family, that inertia
persists throughout the interval. The polynomial continuation
at nonintegral $k$ need not represent a graph.

There are exactly 42 exceptional integers before the thresholds:
$k=3,\ldots,7$ for $s=4$ and $k=3,\ldots,39$ for $s=7$.
Each is certified by rational congruence, as are all eighteen
$k=2$ cases. The table therefore covers every $k\ge2$ for
each $2\le s\le19$. The size never exceeds eight, and is five
throughout $9\le s\le19$.

These exceptional block sizes are substantive.
For $s=7$ and $6\le k\le24$, the initial seven-profile block has
inertia $(1,6,0)$, whereas the initial eight-profile block has
inertia $(2,6,0)$. For $s=4$ and $k=3,4,5$, the five-profile
block also fails and the six-profile block works.
This does not prove a lower bound on the number of target
states for other constructions.

\begin{proof}[Proof of Theorem~\ref{thm:main}]
If $m$ is even, deleting an edge of $C_m$ leaves its independence
number unchanged. The six-state theorem of \cite{SixSpin} applies.
For odd $m=2k+1$, Lemma~\ref{lem:ground} and
\eqref{eq:occupation} give the exact leading kernel.
Proposition~\ref{prop:large} and the bounded-size certificates
produce a nonsingular initial block with at least two positive
eigenvalues, of size at most eight for every $s\ge2$ and of size
five for every $s\ge9$.
Lemma~\ref{lem:compression} supplies at most ten and seven states,
respectively.
\end{proof}

\section{An exact positive target and reproducible evidence}
For $C_5[K_9]$, take seven nodes
\begin{equation}
 a=\left(-\frac{4095}{4096},\,0,\,
 -\frac{15}{16},-\frac78,-\frac{13}{16},
 -\frac34,-\frac{11}{16}\right).                       \label{eq:witness}
\end{equation}
The target $K=\one\one^\top-aa^\top$ is strictly positive
and has inertia $(1,1,5)$.
For $h=17$, the restriction to the five auxiliary variables of
\[
 \Hess(\partial_0^{17}\partial_1^{26}Z_{C_5[K_9]})
\]
has inertia $(2,3,0)$.
This is a full finite-parameter computation using all 4335
occupation vectors of total size 17, rather than just the
leading layer. The supplementary certificate retains the
rational target, scaled integer Hessian, positive normalization
and rational elimination pivots. It is possible to check its
inertia without any floating-point spectral calculation.

The pin value matters. With the same auxiliary nodes, the
earlier trial $\ell=-255/256$ gives auxiliary inertia $(1,4,0)$
and fails to certify the desired obstruction. That failed
trial is retained. Likewise, a prior noninitial four-profile
choice fails at $(k,s)=(3,6)$; the initial seven-profile block
used here succeeds. Neither failed trial is silently discarded
from the reproducibility record.

The supplement contains the classified-pattern kernel, symbolic
five-profile construction, large-size coefficient certificates,
bounded-size tail and exception certificates, exact inertia
routine and an executable verifier. The verifier checks identities
and sign coefficients rather than trusting stored PASS labels.
Source-audit reports and their independently written code are
kept separately from the author computations. The source audit
reconstructed the six large-size certificates, their 836
positive coefficients, all 42 exceptional integers and the
positive target \eqref{eq:witness}. Its further finite
cross-checks corroborate the classification; the proof of the
infinite parameter ranges is the symbolic argument above.

All exact certificates are supplementary parts of this argument.
They do not amount to a proof-assistant formalization.
The finite census of other bases from earlier work is not used
to extrapolate the present theorem. The original failed trials
and separate audit outputs remain preserved.

\section{Limitations and further questions}
The theorem concerns uniform cycle clique blowups.
It does not establish the general connected-source conjecture,
unequal clique sizes, an optimal state bound, a fixed target for
the entire family, or a useful uniform denominator bound.
The restriction $s\ge2$ is the domain of the present occupation
argument; no statement about $s=1$ is inferred from its formulas.

The proof separates two difficulties: finding an initial
profile block with two positive eigenvalues and realizing it
with a bounded number of spins. For odd cycles, the total number
of possible profiles grows with $s$, but five initial profiles
suffice once $s\ge9$. An arbitrary later principal minor would
not alone justify this state bound.
The eight-profile exceptions at $s=7$ suggest a specific
optimization problem, but their inertia does not exclude
more efficient auxiliary-node choices.

\paragraph{Assistance and assessment.}
OpenAI Codex assisted with the proofs, writing and exact computation.
The author operates AIRR.
The source argument received a GPT-6 Astra technical audit,
requested at effort \texttt{xhigh}; it found no unresolved
material mathematical objection within the stated scope.
That source audit is not an exact-version manuscript score.
Any subsequent assessment must identify the final PDF hash.
Prior model participation is disclosed; runtime identity and
technical context isolation are not independently attested.
No human peer review, formal verification or worldwide-priority
certificate is claimed.

\begin{thebibliography}{9}
\bibitem{BH}
P. Br\"and\'en and J. Huh.
\emph{Lorentzian polynomials}.
Annals of Mathematics \textbf{192} (2020), 821--891.
\href{https://arxiv.org/abs/1902.03719}{arXiv:1902.03719}.
\bibitem{LOS}
J. Lee, J. Oh and J. Seo.
\emph{Counting homomorphisms in antiferromagnetic graphs via
Lorentzian polynomials}.
\href{https://arxiv.org/abs/2506.13659}{arXiv:2506.13659}, version 2.
\bibitem{LS}
J. Lee and J. Seo.
\emph{Antiferromagnetic models are clique-minimizing}.
\href{https://arxiv.org/abs/2608.17920}{arXiv:2608.17920}, version 1.
\bibitem{RankTwo}
L. Eriksson.
\emph{Rank-two obstructions to Lorentzian graph partition functions}.
AIRR, ARR-2026-7993ZVQWNC9E38P2, version 1, 17 October 2026.
\href{https://airr.science/papers/ARR-2026-7993ZVQWNC9E38P2/versions/v1/}
{Permanent record}. Working paper.
\bibitem{SixSpin}
L. Eriksson.
\emph{Six-spin Lorentzian obstructions for clique blowups of
independence-stable graphs}.
AIRR, ARR-2026-0VEP5BTYB99VNSC0, version 1, 17 October 2026.
\href{https://airr.science/papers/ARR-2026-0VEP5BTYB99VNSC0/versions/v1/}
{Version record}. Working paper.
\bibitem{Turan}
P. Tur\'an.
\emph{On an extremal problem in graph theory} (in Hungarian).
Matematikai \'es Fizikai Lapok \textbf{48} (1941), 436--452.
\end{thebibliography}
\end{document}
